\documentclass[12pt]{article}
\usepackage[utf8]{inputenc}
\usepackage[T1]{fontenc}
\usepackage{hyperref}
\title{Zyphe Chin Language Resources for K-12 Curriculum Development}
\author{Compiled by Claude Sonnet 5.5 (claude-sonnet-5-5) for Andr\'as Kornai}
\date{}
\begin{document}
\maketitle

\section{Introduction}
Zyphe (also spelled Zophei; ISO 639-3 zyp) is a Kuki-Chin language of the Sino-Tibetan family. It is spoken mainly in Thantlang township, Chin State, Myanmar, with some speakers in India. Estimates put it at roughly 17,000 speakers in Myanmar and 3,000 in India. A Roman-based orthography was created in 1998. The search found no school textbooks or primers in Zyphe. Little material exists beyond linguistic documentation and religious media.

\section{Primary Educational Resources}
The search found no primers, textbooks or school programmes for Zyphe. This section is empty because nothing was found.

\section{Linguistic and Reference Materials}
\begin{itemize}
\item Zophei Language Resource (UNT Digital Library): a collection with a PhD thesis on Zophei sound systems, a word-list spreadsheet and audio. \url{https://digital.library.unt.edu/explore/collections/ZYPLR/}
\item The Sound Systems of Zophei Dialects and Other Maraic Languages (UNT Digital Library). \url{https://digital.library.unt.edu/ark:/67531/metadc1913082/}
\item ASJP Chin Zyphe wordlist. \url{https://asjp.clld.org/languages/CHIN_ZYPHE}
\item Wikipedia article on the Zyphe language. \url{https://en.wikipedia.org/wiki/Zyphe_language}
\end{itemize}

\section{Religious and Community Resources}
\begin{itemize}
\item Scripture Earth page for Zyphe, which lists Scripture resources. \url{https://www.scriptureearth.org/00eng.php?iso=zyp}
\item Joshua Project profile of the Zophei people in Myanmar. \url{https://joshuaproject.net/people_groups/19593/BM}
\end{itemize}

\section{Digital and Audio Resources}
\begin{itemize}
\item Global Recordings Network, Zyphe: audio and video recordings. \url{https://globalrecordings.net/en/language/zyp}
\end{itemize}

\section{Curriculum Development Context}
Zyphe is a small, mostly oral language with a short written tradition. It is not known to be taught in schools. A curriculum would have to start from the documentation above (word lists, phonology, Scripture texts) and from community work on the orthography. Every school-level resource would need to be newly created.

\end{document}